% ============================================================================
% Chapter 4 — Entity Component System
% ============================================================================
\chapter{Entity Component System}

The ECS is the backbone of every Caffeine game. All game objects are
\emph{entities}; all data lives in \emph{components}; all logic runs in
functions that iterate entities with specific sets of components.

The implementation uses an \emph{archetype} layout: entities that share
exactly the same set of component types are stored together in contiguous
memory. Iteration is therefore cache-friendly regardless of entity count.

\section{Entity}

\textbf{Header:} \texttt{ecs/Entity.hpp}

An entity is a lightweight handle — a numeric ID plus a pointer back to its
\texttt{World}. Entities are created by the world and should be stored by
value.

\begin{lstlisting}
Entity e = world.create("Bullet");
bool alive = e.isValid();
\end{lstlisting}

Entities expose shorthand methods that forward to the world:

\begin{longtable}{ll}
\toprule
\textbf{Method} & \textbf{Description} \\
\midrule
\texttt{e.add<T>(args...)}  & Add component T (returns \texttt{T\&}) \\
\texttt{e.remove<T>()}      & Remove component T \\
\texttt{e.has<T>()}         & Returns \texttt{bool} \\
\texttt{e.get<T>()}         & Returns \texttt{T*} (null if absent) \\
\texttt{e.getOrAdd<T>()}    & Returns \texttt{T\&}, adding if needed \\
\texttt{e.isValid()}        & \texttt{bool} — false for destroyed entities \\
\bottomrule
\end{longtable}

\section{World}

\textbf{Header:} \texttt{ecs/World.hpp}

\texttt{World} owns all entities and their component data. Typically one
world exists per scene.

\begin{lstlisting}
World world;
\end{lstlisting}

\subsection{Entity lifecycle}

\begin{longtable}{ll}
\toprule
\textbf{Method} & \textbf{Description} \\
\midrule
\texttt{world.create(name)}   & Create entity; \texttt{name} is optional debug label \\
\texttt{world.destroy(e)}     & Destroy entity and all its components \\
\texttt{world.destroyAll()}   & Destroy every entity in the world \\
\texttt{world.entityCount()}  & \texttt{u32} total live entities \\
\bottomrule
\end{longtable}

\subsection{Component management}

\begin{lstlisting}
// Add a component (constructed in-place)
Transform& t = world.add<Transform>(e);

// Add with constructor arguments
Velocity2D& v = world.add<Velocity2D>(e, 0.0f, -9.8f);

// Check / get
if (world.has<Health>(e)) {
    Health* hp = world.get<Health>(e);   // non-owning pointer
    hp->current -= 10;
}

// Remove
world.remove<Tag>(e);
\end{lstlisting}

\subsection{Iterating entities — forEach}

\texttt{forEach} visits all entities that have \emph{at least} the listed
component types. The callback receives the entity plus references to the
requested components.

\begin{lstlisting}
#include "ecs/ComponentQuery.hpp"

ComponentQuery q;
q.require<Transform>().require<Velocity2D>();

world.forEach<Transform, Velocity2D>(q,
    [&](Entity e, Transform& t, Velocity2D& v) {
        t.position.x += v.x * dt;
        t.position.y += v.y * dt;
    }
);
\end{lstlisting}

\textbf{Important:} do not add or remove components while inside a
\texttt{forEach} callback — this changes archetype membership and invalidates
the internal iterators.

\subsection{Parallel iteration — forEachParallel}

For CPU-heavy systems (physics integration, pathfinding), dispatch the
iteration across the engine's job system:

\begin{lstlisting}
world.forEachParallel<Transform, Velocity2D>(q, jobSystem,
    [](Entity e, Transform& t, Velocity2D& v) {
        t.position.x += v.x * dt;
        t.position.y += v.y * dt;
    }
);
// forEachParallel blocks until all jobs complete.
\end{lstlisting}

The engine splits entities into batches of 1000 and dispatches one job per
batch. Callbacks must be thread-safe with respect to each other.

\section{ComponentQuery}

\textbf{Header:} \texttt{ecs/ComponentQuery.hpp}

A query describes which component types are required (or excluded). Build one
once and reuse it each frame.

\begin{lstlisting}
ComponentQuery query;
query.require<Transform>()
     .require<Velocity2D>()
     .require<Sprite>();
\end{lstlisting}

Pass the same query to \texttt{world.forEach} or \texttt{world.forEachParallel}.

\section{Typical Pattern — System as a Free Function}

\begin{lstlisting}
void movementSystem(World& world, f32 dt) {
    static ComponentQuery q = []() {
        ComponentQuery q;
        q.require<Transform>().require<Velocity2D>();
        return q;
    }();

    world.forEach<Transform, Velocity2D>(q,
        [dt](Entity, Transform& t, Velocity2D& v) {
            t.position.x += v.x * dt;
            t.position.y += v.y * dt;
        }
    );
}
\end{lstlisting}

The query is constructed once (static local) and reused every frame, which
avoids the small allocation cost of rebuilding the component-set bitmask.
